\newpage
\setlength{\absparsep}{18pt} % ajusta o espaçamento dos parágrafos do resumo
\begin{resumo}[Resumo]
% ---
% O trabalho apresentado surge como uma proposta de projeto conceitual para a solução de problemas de locomoção que ocorrem dentro da Universidade de Brasília. A partir de levantamentos sobre o tamanho dos Campi da universidade e do distanciamento entre estes, observou-se que as longas distâncias são fatores dificultadores da experiência cotidiana da comunidade universitária. Desta forma, optou-se por um sistema de transporte multimodal utilizando patinetes elétricos (e-scooter) para o transporte intracampi integrados ao serviço intercampi já oferecido pela instituição. A partir desta hipótese de resolução são apresentados, neste trabalho, requisitos, entregáveis, indicadores, entre outros aspectos relacionados à possibilidade de sua implantação. 
O trabalho apresentado surge como uma proposta de projeto conceitual para a solução de problemas de locomoção que ocorrem dentro da Universidade de Brasília. A partir de levantamentos sobre o tamanho dos Campi da universidade e do distanciamento entre estes, observou-se que as longas distâncias são fatores dificultadores da experiência cotidiana da comunidade universitária. Desta forma, optou-se por um sistema de transporte multimodal utilizando patinetes elétricos (\estrangeirismo{e-scooter}) para o transporte intracampi integrados ao serviço intercampi já oferecido pela instituição. A partir desta hipótese de resolução são apresentados, neste trabalho, requisitos, entregáveis, indicadores, entre outros aspectos relacionados à possibilidade de sua implantação. 
% ---
\vspace{\onelineskip}

\noindent
\textbf{Palavras-chave}: \pchaveinome. \pchaveiinome. \pchaveiiinome. \ifthenelse{\equal{\pchaveivnome}{}}{}{\pchaveivnome.}
\end{resumo}